% No preamble. Prefix abw-. Audited research, 17 October 2026.
\section{Uniform actual-base heat comparison from nonlinear Ward jets}
\label{abw-section}

This section concerns every applicable joint infinite structural system
with the established homogeneous classical inputs below. Constants are
uniform in the old classical depth \(M\). We do not interpolate a finite
momentum grid or identify finite winding kernels with periodizations.

Let \(S_M\) be the actual homogeneous classical minimum action in current
unit coordinates, with a fixed \(4\dim\mathfrak g\) input fiber.
Use its translation, hypercubic, reality and nonlinear gauge Ward
identities, and its uniformly exponentially rooted kernels through field
order four. The needed finite spatial moments, local exponential-chart
derivatives and their constants are uniform in \(M\).
Write \(S_W=S_1\) for the normalized Wilson action. The established flat
normalization and lower bound are
\begin{equation}
 q_M\ge q_W=d_{\rm curl}^{T}d_{\rm curl},\qquad q_MD=0,\qquad
 q_{M,\mathrm{mid}}(p)=|p|^2I-pp^T+O(|p|^4).
 \label{abw-input}
\end{equation}
These are the actual classical current-unit bounds. A bound for an
arbitrary quantum activity with constants depending on a newly chosen
auxiliary grouping would not replace them.

The background-Feynman vector base, in the actual left-relative link
chart, is
\[
 {\cal B}_M(a)=S_M''(\exp(ia))+D_aD_a^T.
\]
For \(\tau(H^2)=1\) and distinct directions \(\alpha,\nu\), define its
affine heat coefficient by its convergent fixed-time Taylor kernels:
\begin{equation}
 h_M(s)=[t^2]\operatorname{tr}_{E}
 \bigl(e^{-s{\cal B}_M(a(t))}\bigr)_{00},\qquad
 a_\alpha(t,x)=t x_\nu H.
 \label{abw-h}
\end{equation}
This is an infinite-lattice density, not a varying-flux family on a
finite torus.

\begin{theorem}[Uniform base comparison]
\label{abw-main}
Under these inputs there is \(C\), independent of \(M\), such that
\begin{equation}
 |h_M(s)-h_W(s)|\le C s^{-1}\quad(s\ge1),\qquad
 \int_1^{K^2}|h_M(s)-h_W(s)|\,\frac{ds}{s}\le C.
 \label{abw-bound}
\end{equation}
The scalar base \(D_a^TD_a\) is identical for \(M\) and \(W\).
Thus the same comparison applies to the vector-minus-two-scalar
combination. Explicit periodizations may also be used where their
identification has been established.
\end{theorem}

\subsection{Actual relative vertices, including chart curvature}
\label{abw-chart-section}

Set \(R_M=S_M-S_W\). Its Maxwell coefficient vanishes by
\eqref{abw-input}. The previously proved total nonlinear Ward extraction
therefore gives, with uniform analytic moment bounds,
\begin{equation}
 R_2=O(|p|^4),\qquad R_3=O((\max_i|p_i|)^3),\qquad
 R_4=O((\max_i|p_i|)^2).
 \label{abw-ward-orders}
\end{equation}
Analytic changes between midpoint and source-site bond phases preserve
these orders.

For \(\Xi_a(B)=\log(e^{iB}e^{ia})/i\), write
\[
 J(a)=D_B\Xi_a(0)=I+J_1(a)+J_2(a,a)+O(a^3),\qquad
 T(a)=D_B^2\Xi_a(0)=T_1(a)+O(a^2).
\]
In particular \(T(0)=0\). The exact relative Hessian is
\begin{equation}
 L_R(a)[B,C]=R_M''(a)[J(a)B,J(a)C]
                  +R_M'(a)[T(a)[B,C]].
 \label{abw-chain}
\end{equation}
At background degree one it is
\[
 R_3[a,B,C]+R_2[J_1(a)B,C]+R_2[B,J_1(a)C].
\]
Its degree-two coefficient is
\begin{align}
 &\tfrac12R_4[a,a,B,C]
 +R_3[a,J_1(a)B,C]+R_3[a,B,J_1(a)C]\notag\\
 &\quad+R_2[J_2(a,a)B,C]+R_2[B,J_2(a,a)C]
 +R_2[J_1(a)B,J_1(a)C]+R_2[a,T_1(a)[B,C]].
 \label{abw-second-chart}
\end{align}
Mixed derivatives are obtained by polarization. The last term is the
current times chart-curvature term and is retained.

All coefficients of \(J_1,J_2,T_1\) are local at one bond with fixed color
norms. They only merge momenta and introduce no inverse momentum.
Thus the flat, first-background and second-background differences of the
actual Hessians have respective momentum orders \(4,3,2\).
The gauge-fixing term \(D_aD_a^T\) is the same for both actions and
cancels from these differences.

\subsection{Fourier bounds and exact total heat variations}
\label{abw-diagram-section}

Put \(B_M(k)={\cal B}_M(0,k)\) and
\(\delta(k)=4\sum_\mu\sin^2(k_\mu/2)\). Then
\begin{equation}
 B_W(k)=\delta(k)I,\qquad
 c\delta(k)I\le B_M(k)\le C\delta(k)I.
 \label{abw-ellipticity}
\end{equation}
The upper bound follows from the uniform quadratic zero and bounded
second moments. On a fixed small punctured ball, with \(\rho=|k|\),
\begin{equation}
 \|\partial^r(B_M-B_W)\|\le C\rho^{4-r}\ (r\le2),\qquad
 \|\partial B_M\|\le C\rho,\qquad \|\partial^2B_M\|\le C.
 \label{abw-flat-jets}
\end{equation}

Let \(V_a^M(k,p)\) be the first vertex for a background
\(e^{ipx}e_a\), mapping quantum momentum \(k\) to \(k+p\).
Let \(W_{ab}^M(k;p,-p)\) be the mixed second vertex, preserving \(k\).
For \(\Delta=M-W\), \eqref{abw-ward-orders}--\eqref{abw-second-chart}
imply at \(p=0\), for \(r\le2\),
\begin{align}
 \|\partial_p^r\Delta V_a(k,p)\|
 +\|\partial_p^r\Delta V_b(k+p,-p)\|
       &\le C\rho^{3-r},\notag\\
 \|\partial_p^r\Delta W_{ab}(k;p,-p)\|
       &\le C\rho^{2-r}.
 \label{abw-vertex-difference}
\end{align}
The second first-vertex term includes differentiation of both arguments.
For either baseline it suffices to use
\begin{equation}
 \|\partial_p^rV\|\le C\rho^{1-r},\qquad
 \|\partial_p^rW\|\le C\rho^{-r}.
 \label{abw-baseline-vertices}
\end{equation}
Negative powers here are only loose bounds on analytic functions.
In particular \(\partial_p^2W\) is bounded. Away from the small ball all
required vertices and derivatives are uniformly bounded.

For compact source variations the second total heat variation is
well-defined by trace-class Duhamel insertions. There is no need to assign
a value to the extensive vacuum trace. Differentiating
\(D\operatorname{Tr}e^{-sB}[a]=-s\operatorname{Tr}(e^{-sB}B'[a])\)
gives
\begin{equation}
 D^2\operatorname{Tr}e^{-sB}[a,b]
 =-s\operatorname{Tr}(e^{-sB}B''[a,b])
 +s^2\int_0^1\operatorname{Tr}
  (e^{-s(1-v)B}B'[a]e^{-svB}B'[b])\,dv.
 \label{abw-total-duhamel}
\end{equation}
The integral already includes both orders; there is no further factor
two. Only the \emph{total} trace was cycled.

The external-momentum kernel is accordingly
\begin{align}
 {\cal H}_{ab}^M(s,p)
 &=-s\int_{\mathrm{BZ}}\operatorname{tr}
       (e^{-sB_M(k)}W_{ab}^M(k;p,-p))\,\frac{dk}{(2\pi)^4}
       \notag\\
 &\quad+s^2\int_0^1\int_{\mathrm{BZ}}\operatorname{tr}
 \left(e^{-s(1-v)B_M(k+p)}V_a^M(k,p)
       e^{-svB_M(k)}V_b^M(k+p,-p)\right)
                          \frac{dk}{(2\pi)^4}\,dv.
 \label{abw-diagram}
\end{align}
At fixed \(s\) the kernels have the moments permitting two external
derivatives under this integral.

\subsection{Endpoint-safe heat estimates and the gain}
\label{abw-estimate-section}

Duhamel differentiation and \eqref{abw-ellipticity} give, for
\(E_M(t,k)=e^{-tB_M(k)}\) and \(r=0,1,2\),
\begin{equation}
 \|\partial_k^rE_M(t,k)\|
 \le C_r\rho^{-r}(1+t\rho^2)^r e^{-ct\rho^2}.
 \label{abw-heat}
\end{equation}
The sharper first and second derivative factors are respectively
\(Ct\rho\) and \(C(t+t^2\rho^2)\). No spectral frame is differentiated.
The exact difference identity
\[
 E_M(t)-E_W(t)
 =-\int_0^t E_M(t-u)(B_M-B_W)E_W(u)\,du
\]
and \eqref{abw-flat-jets} give
\begin{equation}
 \|\partial_k^r(E_M-E_W)(t,k)\|
 \le C_r t\rho^{4-r}(1+t\rho^2)^r e^{-ct\rho^2}.
 \label{abw-heat-difference}
\end{equation}
These bounds hold for every \(t\ge0\); in particular they contain no
singularity at either endpoint \(v=0,1\) of \eqref{abw-diagram}.

Differentiate \eqref{abw-diagram} twice externally and subtract.
For the tadpole, the two differences are
\((E_M-E_W)\partial_p^2W_W\) and
\(E_M\partial_p^2(W_M-W_W)\).
Before the outer factor \(s\), they are bounded by
\(Cs\rho^4e^{-cs\rho^2}\) and \(Ce^{-cs\rho^2}\).
Their integrals are \(O(s^{-2})+O(s^{-1})\) after that factor.

Telescope the bubble product into its four factor differences.
Replacing a vertex gives combined vertex degree \(3+1=4\).
After distributing two derivatives between vertices and heat factors,
the full product is bounded by
\[
 C\rho^2 P(s\rho^2)e^{-cs\rho^2},
\]
where \(P\) is a fixed polynomial. Replacing a heat factor instead
gives, by \eqref{abw-heat-difference}, a time at most \(s\) and degree
\(4+1+1\); the differentiated bound is
\[
 Cs\rho^4 P(s\rho^2)e^{-cs\rho^2}.
\]
Both bounds are uniform in \(v\), since the two heat times sum to \(s\).
The four-dimensional integrals of these expressions are
\(O(s^{-3})\) and \(O(s^{-3})\), respectively.
The prefactor \(s^2\) therefore gives \(O(s^{-1})\).
Away from zero, a fixed spectral gap gives a polynomial in \(s\)
times \(e^{-cs}\). We have proved
\begin{equation}
 |\partial_p^2({\cal H}_{ab}^M-{\cal H}_{ab}^W)(s,0)|
 \le C/s,\qquad s\ge1.
 \label{abw-two-point-bound}
\end{equation}

\subsection{Density normalization and scope}
\label{abw-density-section}

The diagonal color heat density is gauge invariant, and its linear
color jet is zero for a semisimple group. The established density
Ward/reanchoring identity at fixed \(s\) consequently gives
\begin{equation}
 h_M(s)=\tfrac14\,
 \partial_{p_\nu}^2{\cal H}_{\alpha H,\alpha H}^M(s,0).
 \label{abw-density-factor}
\end{equation}
Indeed the Maxwell extraction is both twice the affine quadratic
density and one half of this total Hessian derivative. On the diagonal
component the source-site and midpoint phases cancel.

For the actual endpoint-\(K\) source \(a_\alpha=t x_\nu H/K^2\),
translation changes the background only by a constant pure gauge.
Every diagonal coefficient in a final block is therefore the same.
The \(K^4\) sites cancel the \(K^{-4}\) from the two source factors.
Its block coefficient is exactly \eqref{abw-h}; no partial-trace
cyclicity or finite-torus flux interpolation enters.
Equations \eqref{abw-two-point-bound} and \eqref{abw-density-factor}
prove Theorem~\ref{abw-main}.

This result supplies the uniform old-depth comparison for the base
operator. The soft-penalty, infrared, short-time and coarse-determinant
comparisons remain separate inputs to an actual determinant assembly.
The proof does not by itself construct a finite-regulator running
trajectory, a large-field RG history, a continuum measure or a mass gap.

